\section{Discussion}
\label{sec:discussion}

Both SAC and TD3 proved to be well-suited algorithms for the hockey environment. SAC 
converged faster and required less manual tuning due to its automatic entropy 
regularization, while TD3 with pink noise exploration provided a competitive alternative. 
Hyperparameter tuning required considerable effort, but the most significant performance 
gains ultimately came from longer training runs and exposure to diverse opponents rather 
than from any single algorithmic choice.

The modular YAML-based framework proved valuable throughout the project, allowing rapid 
iteration over configurations without code changes. Combined with TensorBoard logging, 
we maintained a clear picture of training progression at all times. The round-robin 
tournament with TrueSkill ranking provided a reliable and interpretable comparison 
across all agent variants at every stage of development.

\medskip
The most impactful finding was that training against a diverse pre-trained opponent pool 
led to substantially stronger and more generalizable agents. Agents trained only against 
the basic opponents or against pure self-play snapshots tended to overfit quickly, 
converging to narrow strategies that exploited specific weaknesses of familiar opponents. 
This overfitting became apparent only through the tournament, where top-ranked agents 
unexpectedly lost to stylistically different lower-ranked opponents. The final SAC and 
TD3 agents trained against the pre-trained pool win consistently against both the weak 
and strong basic opponent, as reflected in the win rates reported in 
Table~\ref{tab:tournament_results} in the appendix. In hindsight, we would have invested 
more time earlier in building a larger and more varied pool, including agents trained 
under more diverse configurations and potentially additional algorithm types, to capture 
an even broader range of play styles.

Our results strongly suggest that exposure to diverse opponents is the primary driver of 
generalization. For the live tournament against the other course participants, fine-tuning 
against recorded opponent trajectories between matches could therefore be a particularly 
effective strategy to gain ranking, as it would allow the agent to adapt directly to the 
playing styles encountered in the competition.

\medskip
Overall, the project offered valuable practical insight into the challenges of 
reinforcement learning that goes well beyond theoretical study alone. The hands-on 
nature of iteratively training, analyzing games, and refining strategies made the 
underlying complexity of the problem genuinely concrete. The infrastructure built 
around the training pipeline provided a solid foundation that made this iterative 
process efficient, reproducible, and enjoyable.